\section{Einleitung}\label{sec:einleitung}

Zu Beginn dieser Vorlage möchten wir betonen, dass nicht alle Teile der Vorlage für jede Abschlussarbeit relevant sind. Wenn es keine Tabellen in ihrer Arbeit gibt, ist es selbstverständlich nicht notwendig ein Tabellenverzeichnis anzulegen. Diese Abschnitte dienen nur der einfachen Verwendungen, falls ein solcher Abschnitt benötigt wird. Andere Abschnitte wie die Einleitung sind dagegen unumgänglich.\\

Die Einleitung führt in das Thema ein und erläutert die \textbf{Fragestellung} Ihrer Arbeit sowie das \textbf{methodische Vorgehen}. Die Einleitung sollte insgesamt die folgenden Fragen beantworten: 
\begin{itemize}
    \item Worum geht es in der Arbeit?
    \item Warum ist das Thema wichtig?
    \item Was möchten Sie herausfinden bzw. was ist das Problem?
    \item Welche wissenschaftliche Methodik wird angewandt?
    \item Welche Forschungsfragen werden beantwortet?
\end{itemize}

\textbf{\uline{Hinweise}}
\begin{itemize}
    \item Mit der Einleitung beginnt die Nummerierung der Seiten mit arabischen Nummern (bis zu letzten Seite des Fazits).
    \item Es empfiehlt sich die Einleitung, zusammen mit dem Fazit, erst zum Ende der Ausarbeitung der Arbeit zu verfassen.
\end{itemize}

\textbf{\uline{Detailliertere Anregungen zur Einleitung}}
Stellen Sie in der Einleitung einen Bezug Ihrer Arbeit zu allgemein diskutierten \textbf{aktuellen Themen und übergeordneten Fragestellungen} her. Damit leiten Sie die Bewegründe und \textbf{Relevanz} Ihrer spezifischen Forschungsfrage bzw. Problemstellung im Kontext des Forschungsumfelds ab.\\
\\
\textbf{\uline{Hinweise}}
\\
Die Schilderung der Ausgangssituation bzw. Problemstellung ist mit \textbf{Quellen zu belegen!}

Die Zielstellung \textbf{konkretisiert die Forschungsfrage}, die im Rahmen der Arbeit beantwortet wird und zeigt die zu \textbf{erwartenden Ergebnisse auf}. Die in der Regel notwendige Abgrenzung bzw. Eingrenzung des Themas sollte sinnvoll begründet werden. 

Abschließend wird die Vorgehensweise zur Erreichung des Forschungsergebnisses an dieser Stelle skizziert, d.~h. es ist ein knapper Überblick über die anschließenden \textbf{Kapitelinhalte} zu geben. Dazu wird dieser kurz für jedes der folgenden Kapitel dargestellt, sodass der Leser einen Eindruck von der Struktur der Arbeit bekommt.